\documentclass[8pt]{beamer}
\usepackage{tikz}
\usepackage{tikz-cd}
\usepackage{qtree}
\usepackage{amsfonts}
\usepackage{amssymb}
\usepackage{amsmath}
\usepackage{amscd}
\usepackage{amsthm}
% \usepackage{mathpartir}
\usepackage{bbm}
\usepackage[mathcal]{euscript}
% \usepackage[utf8]{inputenc}
\setbeamertemplate{navigation symbols}{}%remove navigation symbols
%\usepackage[T1]{fontenc}
\usepackage{lmodern} \title{Are Rindler quanta real?}
\author{Hans Halvorson}
\date[November 2013]{November 12, 2013} \usetheme{Copenhagen}
\begin{document}
\def\2{\mathcal}
\def\7{\mathbb}
\newenvironment{proenv}{\only{\setbeamercolor{local structure}{fg=green}}}{}
\newenvironment{conenv}{\only{\setbeamercolor{local structure}{fg=red}}}{}  
\newenvironment{neuenv}{\only{\setbeamercolor{local
      structure}{fg=yellow}}}{}


\newcommand{\3}{\mathfrak}

\begin{frame}[plain]
  \titlepage
\end{frame}

\begin{frame}{first quantization}

  Begin with a classical phase space $S$, a real vector space equipped
  with a sympletic form $\sigma$, and a one-parameter group $t\mapsto
  T_t$ of symplectomorphisms.
  
  \bigskip \textbf{first quantization} is equivalently:

\begin{itemize}
\item a linear embedding $\eta$ of $S$ into a Hilbert space $\2H$ such
  that $\eta (T_{t}f) = U_{t}\eta (f)$.
\item a complex structure $J$ on $S$.
\item a real inner product $\mu$ on $S$.
\end{itemize}


\end{frame}
\begin{frame}{second quantization}

  \bigskip \textbf{second quantization} takes the one-particle space
  $\2H$ to Fock space
$$ \2F (\2H ) \: = \: \7C\oplus \2H\oplus (\2H\otimes \2H)\oplus \cdots .$$
Observables and symmetries on $\2H$ also lift to observables and
symmetries on $\2F (\2H )$.

\bigskip \textbf{vacuum state} $\Omega = 1\oplus 0\oplus 0\oplus
\cdots $

\end{frame}

\begin{frame}{creation and annihilation operators}
\[
\begin{array}{lll} a_n^*(f)(f_1\otimes \cdots \otimes f_{n-1}) & = & f\otimes f_1\otimes \cdots \otimes f_{n-1} \\ \\ 
  a_n(f)(f_1\otimes\cdots\otimes f_{n-1}) & = & \langle f,f_1\rangle
  _J f_2\otimes\cdots \otimes f_{n-1} \end{array} \]

\bigskip

\[
\begin{array}{lll}
  a^*(f) & = & a_1^*(f) \oplus \sqrt{2}a^*_2(f)\oplus \sqrt{3}a_3^*(f)\oplus \dots \\ \\
  a(f) & = & 0\oplus a_1(f)\oplus \sqrt{2}a_2(f)\oplus \sqrt{3}a_3(f)\oplus \cdots \end{array} \]




\end{frame}
\begin{frame}{number and field operators}

\[ 
\begin{array}{lll}
N(f) & = & a^*(f)a(f) \\ \\
N & = & \sum _{i=1}^\infty N(f_i) \\ \\ 
\Phi (f) & = & \frac{1}{\sqrt{2}}(a^*(f)+a(f)) \\ \\
W(tf) & = & \exp (it\Phi (f)) \\ \\ 
\langle \Omega ,W(f)\Omega\rangle & = & e^{-\langle f,f\rangle _J/4} \end{array} \]


\end{frame}


\begin{frame}{the Weyl algebra}

  \textbf{objective:} capture the algebraic structure of the canonical
  commutation relations
$$ Q_nP_m - P_mQ_n = -i\delta _{nm}\hbar I .$$

\bigskip \textbf{step 1 (exponentiate)}
$$ e^{iaQ}e^{ibP} = e^{-iab/2}e^{ibP}e^{iaQ} .$$

\bigskip \textbf{step 2 (collect)} $W(a,b)=e^{iab/2}e^{ibP}e^{iaQ}$
$$ W(x)W(y) = e^{-i\sigma (x,y)}W(x+y). $$

\end{frame}
\begin{frame}{the Weyl algebra}

  \textbf{fact:} There is a unique smallest $C^*$-algebra $\2A
  [S,\sigma ]$ generated by unitary operators $\{ W(x):x\in S \}$
  satisfying the Weyl commutation relations.

  \bigskip Regular representations of the algebra $\2A [S,\sigma ]$
  correspond to representations of the CCR.

\bigskip Fock representations of $\2A [S,\sigma ]$ correspond to:
\begin{itemize}
\item complex structures on $(S,\sigma )$
\item real inner products on $(S,\sigma )$
\item quasifree states on $\2A [S,\sigma ]$
\end{itemize}

\end{frame}

\begin{frame}{inequivalent Fock representations}

  \textbf{fact} In the Minkowski vacuum representation $\pi _M$, the
  quasilocal algebra $\pi _M(\2A [S,\sigma ])''$ is type III.

  \bigskip \textbf{corollary} The Rindler and Minkowski vacuum
  representations are \textcolor{blue}{disjoint}.



\end{frame}


\begin{frame}{Fulling's theorem} 

  Two Fock vacuum representations $(\2H,\Omega )$ and $(\2H',\Omega
  ')$ of the CCRs are unitarily equivalent if and only if:
\[ \langle \Omega ' |N\Omega '\rangle <\infty \] 

\vspace{1em} The Minkowski and Rindler vacuum representations are
inequivalent.  Therefore, \[ \langle \Omega _{m}| N_{r}\Omega
_{m}\rangle = \infty \]

\end{frame}

\begin{frame}{Chaiken's theorem}

  Let $(\2H _m,\Omega _{m})$ be the Minkowski vacuum representation of
  the CCRs.  Then the Rindler number operator $N_{r}$ cannot be
  defined on the Minkowski Hilbert space $\2H _m$.  In fact, there is
  no Hilbert space on which both $N_m$ and $N_r$ can be simultaneously
  defined.

  \bigskip \textbf{corollary.} The expression $N_{m}N_{r}-N_{r}N_{m}$
  doesn't make any sense.

  \bigskip \textbf{corollary.} The expression $\langle \Omega
  _{m}|N_{r}\Omega _{m}\rangle =\infty \,$ also doesn't make sense!

\end{frame}

\begin{frame}{ }

  \begin{itemize}
\item A quantity cannot take an infinite value.
\item $N_r$ cannot be applied to the vector $\Omega _m$.
\end{itemize}

\end{frame}


\end{document}